\documentclass[10pt,a4paper]{amsart}
\usepackage[margin=19mm]{geometry}
\usepackage{amsmath,amssymb,mathtools}
\usepackage[scr=rsfs,cal=euler]{mathalfa}
\usepackage{xfrac}
\usepackage[unicode,hidelinks,bookmarks=false]{hyperref}
\newcommand{\ie}{i.e.\ }
\newcommand{\cf}{cf.\ }
\newcommand{\resp}{respectively\ }
\newcommand{\Subsection}{\subsection}
\providecommand{\nobreakdash}{\nobreak\hbox{-}}
\setlength{\emergencystretch}{2em}
\multlinegap0pt
\usepackage{dsfont}
\usepackage{amssymb}
\usepackage{xcolor}
\newtheorem{lemma}{Lemma}
\newtheorem{theorem}{Theorem}
\newtheorem{remark}{Remark}
\newcommand{\B}{\mathcal{B}}
\DeclareMathOperator{\D}{Der}
\DeclareMathOperator{\I}{Id}
\DeclareMathOperator{\ad}{ad}
\newcommand{\black}{\color{black}}
\DeclareMathOperator{\e}{\varepsilon}
\title{Differential varieties of upper triangular matrices}
\author{Daniela La Mattina, Carla Rizzo}
\date{Source: arXiv:2608.11032v1 (CC BY 4.0)}
\begin{document}
\maketitle

\noindent\emph{Source and licence.} Differential varieties of upper triangular matrices. Version arXiv:2608.11032v1 (CC BY 4.0), distributed by arXiv under the Creative Commons Attribution 4.0 International licence, http://creativecommons.org/licenses/by/4.0/, as recorded in the arXiv licence metadata for this exact version and shown on the abstract page https://arxiv.org/abs/2608.11032v1. Full licence notice: \texttt{licenses/arxiv-2608.11032-CC-BY-4.0.txt}.\par\smallskip

\noindent\textit{Editorial scope.} Counted unit: the article's theorem e1-e2 (source0.tex lines 1184--1194). The article's complete original proof of it is delivered with it (source0.tex lines 1195--1371, one \texttt{proof} environment of 177 lines). No mathematics is written, repaired or re-derived: the statement and the proof are the article's own lines, in the article's own notation, and no retained line is substituted or re-flowed.  Retained uncounted as the source-local context the proof needs: lines 711--717; lines 730--741; lines 759--765; lines 809--812; lines 816--825; lines 956--996; lines 979--981; lines 1013--1015; lines 1017--1019; lines 1023--1030; lines 1058--1061; lines 1090--1093; lines 1106--1112; lines 1141--1150; lines 1161--1172.  Nothing was omitted from any retained span, so the package carries no omission record.  No retained line contains a citation, so the wrapper carries no bibliography and no bibliography entry.  No retained line contains a figure environment, an image-inclusion command or drawing code, so no image is delivered and the package declares no asset dependency.  Editorial formatting only, in addition: an AMS \texttt{amsart} preamble that restates the article's own declarations and macros used by the retained lines, namely the environments \texttt{lemma}, \texttt{theorem}, \texttt{remark} on separate counters, together with the standard packages those lines need.  The retained environments print in this document as Lemma 1, Theorem 1, Remark 1 in order of appearance, whereas the article prints the same items with the numbering of its own full document.  The complete native source of the article as distributed in the arXiv e-print for this exact version is bundled unchanged as \texttt{sources/207381/source0.tex}.
\begin{lemma}\label{lem: x_1^{u_1}x_2^{u_2}}
    Let $u_1, u_2\in U(L)$. For any $n\geq 2$, the commutator $[x_1^{u_{1}}x_2^{u_{2}}, x_3, \ldots, x_n]$ can be written as a linear combination of $L$-polynomial of the form 
    \begin{align*}
        &[x_{1}^{u_{1}},x_{i_1},\dots,x_{i_{k}}][x_{2}^{u_{2}},x_{j_1},\dots,x_{j_{n-k-2}}],
    \end{align*}
     where $0\leq k \leq n-2$. 
\end{lemma}

\begin{lemma}\label{lem: ordered comm x_1^u}
    Let $u\in U(L)$. For any $n\geq 3$ and $2\leq a \leq n-1$, we have
$$
[x_1^u,x_2,\ldots, x_a,x_{a+1}, \ldots, x_{n}]= [x_1^u,x_2,\ldots, x_{a-1}, x_{a+1},x_{a}, x_{a+2}, \ldots, x_{n}] + \Gamma,
$$
where $\Gamma$ is a linear combination of $L$-polynomials of the form
\begin{align*}
    &[x_1^u, x_{i_1}, \ldots, x_{i_k}][x_a, x_{a+1}, x_{j_1}, \ldots, x_{j_{n-k-3}}],\\
    &[x_a, x_{a+1}, x_{i_1}, \ldots, x_{i_k}][x_1^u, x_{j_1}, \ldots, x_{j_{n-k-3}}],
\end{align*}
where $0\leq k \leq n-3$.
\end{lemma}

\begin{theorem}\label{Teo: Id^L UT_3 ordinarie}
Let $L$ be a Lie algebra and consider the algebra $UT_3$ where $L$ acts trivially (i.e., $\theta=0$). Then 
there exists a basis $\mathcal{B}_L=\{d_i \mid i \geq 1\}$ of $L$ over $F$ such that 
the $T_L$-ideal of $L$-identities of $UT_3$ is generated by the following $L$-polynomials
$$x^{d_i}, \quad  [x_1,x_2][x_3,x_4][x_5,x_6],$$
for all  $i\geq 1.$ Moreover, $c_n^L(UT_3)=(n^2-7n+9) 3^{n-2}+ (3n-6)2^{n-1}+3.$ 
\end{theorem}

\begin{lemma}\label{lem: conseq 3 commutators 1} 
\black
    Let $u\in U(L)$. Then $[x_1,x_2][x_3,x_4][x_5,x_6]\in \langle [x_1,x_2]x_3^{u},\; \big([x_1,x_2]^{u}-[x_1,x_2]\big)[x_3,x_4] \rangle_{T_L}$.
\end{lemma}

\begin{theorem}\label{Teo: Id^L UT_3 epsilon1}
Let $L$ be a Lie algebra and let $UT_3^{\e_1}$  be  the $L$-algebra $UT_3$ where $L$ acts via the one-dimensional Lie subalgebra of $\D(UT_3)$ spanned by the inner derivation $\e_1=\ad_{(-e_{11})}$. Then 
there exists a basis $\mathcal{B}_L=\{d_i \mid i \geq 1\}$ of $L$ over $F$ such that 
the $T_L$-ideal of differential identities of $UT_3^{\e_1}$ is generated by the following $L$-polynomials
$$
x^{d_i}, \quad  x^{d_1^2}-x^{d_1}, \quad [x_1,x_2]x_3^{d_1}, \quad \big([x_1,x_2]^{d_1}-[x_1,x_2]\big)[x_3,x_4],
$$
for all  $i\geq 2.$ Moreover,
 $c_n^L(UT_3^{\e_1})=(n^2-4n)3^{n-2}+ (3n-2)2^{n-1} +2.$ 
\end{theorem}

\begin{remark}[General strategy]\label{rmk: general-strategy}
We record here, once and for all, the five-step scheme used throughout this section to compute
the $T_L$-ideal of identities and the $L$-codimension sequence of $UT_3$ endowed with an
$L$-action; each theorem below only specifies the case-dependent data.

\smallskip
\textbf{Step 1 (adapted basis).} Fix $\theta:L\to\D(UT_3)$ with the prescribed image, and a basis
$\mathcal B_L=\{d_i\mid i\ge1\}$ of $L$ adapted to $\theta$ (i.e., $\theta(d_1)=\delta_1$, or
$\theta(d_1)=\delta_1,\theta(d_2)=\delta_2$ in the two-dimensional case, and $d_i\in\ker\theta$
for the remaining indices).

\smallskip
\textbf{Step 2 (inclusion $I\subseteq\I^L(UT_3^\theta)$).} Let $I$ be the $T_L$-ideal generated by the
polynomials in the statement. Using the explicit action of $\theta(d_1)$ (and $\theta(d_2)$) on
the matrix units, one checks directly that $I\subseteq\I^L(UT_3^\theta)$.

\smallskip
\textbf{Step 3 (PBW reduction).} By the PBW Theorem, every $f\in P_n^L$ is a linear combination of
monomials $x_{i_1}^{u_{a_1}}\cdots x_{i_r}^{u_{a_r}}c_1\cdots c_b$. Case-specific technical lemmas
bound the exponents occurring modulo $I$, the number of commutators ($b\le2$), and force each
non-trivial exponent into a specific position (tail or a prescribed entry of a commutator).  In
particular, for any $u\in \B_{U(L)}$ and $1\le k\le r$, a variable with exponent $u$ occurring in
the tail can be moved into the first entry of a commutator via the identity
\begin{equation}\label{eq: tail-to-commutator}
x_{i_1}\cdots x_{i_{k-1}}x_{i_k}^{u}x_{i_{k+1}}\cdots x_{i_r} = x_{i_1}\cdots x_{i_{k-1}}x_{i_{k+1}}x_{i_k}^{u}x_{i_{k+2}}\cdots x_{i_r} + x_{i_1}\cdots x_{i_{k-1}}[x_{i_k}^{u},x_{i_{k+1}}]x_{i_{k+2}}\cdots x_{i_r},
\end{equation}
applied repeatedly (with $k$ decreasing) to obtain a linear combination of polynomials with the
variable of exponent $u$ in the first entry of a commutator.


\smallskip
\textbf{Step 4 (tail elimination).} Since $UT_3$ has a unit, evaluating the variables of a
maximal-length tail at $1_{UT_3}$ shows, by decreasing induction on the tail length, that every
variable with exponent $1_{U(L)}$ may be assumed to occur inside a commutator.

\smallskip
\textbf{Step 5 (linear independence via evaluations).} For each family of indices indexing a
candidate basis element, a suitable evaluation on matrix units, determined by the action of
$\theta(d_1)$ (and $\theta(d_2)$), isolates its coefficient and shows it must vanish, proving
linear independence; counting the resulting basis gives the $L$-codimension.
\end{remark}

\begin{equation}
x_{i_1}\cdots x_{i_{k-1}}x_{i_k}^{u}x_{i_{k+1}}\cdots x_{i_r} = x_{i_1}\cdots x_{i_{k-1}}x_{i_{k+1}}x_{i_k}^{u}x_{i_{k+2}}\cdots x_{i_r} + x_{i_1}\cdots x_{i_{k-1}}[x_{i_k}^{u},x_{i_{k+1}}]x_{i_{k+2}}\cdots x_{i_r},
\end{equation}

\begin{equation}\label{eq: nu_12}
    (e_{12})^{\nu_{12}}=e_{12}, \quad (e_{11})^{\nu_{12}} = (e_{22})^{\nu_{12}} = (e_{33})^{\nu_{12}}= (e_{23})^{\nu_{12}}=(e_{13})^{\nu_{12}} =0,
\end{equation}

\begin{equation}\label{eq: nu_23}
   (e_{23})^{\nu_{23}}=e_{23}, \quad  (e_{11})^{\nu_{23}} = (e_{22})^{\nu_{23}} = (e_{33})^{\nu_{23}}= (e_{12})^{\nu_{23}}=(e_{13})^{\nu_{23}} =0.
\end{equation}

\begin{lemma}\label{lem: cons IdUT3e1-e2 1}
    Let $d\in L$. If we denote $u_{12}:=2^{-1}(d^2+d)$ and $u_{23}:=2^{-1}(d^2-d)$, then 
    \begin{equation*}
        \begin{split}
          x_1^{u_{12}} x_2^{u_{12}}, \quad x_1^{u_{23}} x_2^{u_{23}} \in \langle x^{d^3}-x^{d}\rangle_{T_L}.
        \end{split}
    \end{equation*}
\end{lemma}

\begin{lemma}\label{lem: cons IdUT3e1-e2 4}
     Let $d\in L$. If we denote $u_{12}:=2^{-1}(d^2+d)$ and $u_{23}:=2^{-1}(d^2-d)$, then 
     $$ x_1^{u_{12}}[x_2,x_3]x_4^{u_{23}} \in \langle x^{d^3}-x^{d}, \ x_1^{u_{12}} \big([x_2,x_3]^{u_{23}}-[x_2,x_3]\big) \rangle_{T_L}.$$
\end{lemma}

\begin{lemma}\label{lem: cons IdUT3e1-e2 3}
    Let $d\in L$. If we denote $u_{12}:=2^{-1}(d^2+d)$ and $u_{23}:=2^{-1}(d^2-d)$, then
    $$ x_1^{u_{12}}x_2 x_3^{u_{12}}, \ x_1^{u_{23}} x_2 x_3^{u_{23}}, \ x_1^{u_{23}}x_2x_3^{u_{12}} \in \langle x^{d^3}-x^{d}, \ x_1^{u_{23}}x_2^{u_{12}}\rangle_{T_L}.$$
\end{lemma}

\begin{lemma}\label{lem: cons IdUT3e1-e2 5}
    Let $d\in L$. If we denote $u_{12}:=2^{-1}(d^2+d)$ and $u_{23}:=2^{-1}(d^2-d)$, then
     $$
 [x_1^{u_{12}},x_2] [x_3,x_4] - [x_2^{u_{12}},x_1] [x_3,x_4]-[x_1,x_2][x_3,x_4]
$$
is a consequence of the $L$-polynomials in $$\bigl\{x^{d^3}-x^{d}, \ x_1^{u_{23}}x_2^{u_{12}}, \ x_1^{u_{23}}[x_2,x_3],\ \big([x_1,x_2]^{u_{12}}-[x_1,x_2]\big)[x_3,x_4]\bigr\}.$$
\end{lemma}

\begin{lemma}\label{lem: cons IdUT3e1-e2 6}
    Let $d\in L$. If we denote $u_{12}:=2^{-1}(d^2+d)$ and $u_{23}:=2^{-1}(d^2-d)$, then
   $$
 [x_1,x_2] [x_3^{u_{23}},x_4] - [x_1,x_2] [x_4^{u_{23}},x_3]-[x_1,x_2][x_3,x_4]
$$
is a consequence of the $L$-polynomials in 
\begin{align*}
   \Bigl\{  x^{d^3}-x^{d}, \ x_1^{u_{23}}x_2^{u_{12}}, \ [x_1,x_2]x_3^{u_{12}}, \ x_1^{u_{12}} \big([x_2,x_3]^{u_{23}}-[x_2,x_3]\big),\ \big([x_1,x_2]^{u_{12}}-[x_1,x_2]\big)[x_3,x_4] \Bigr\}.
\end{align*}
\end{lemma}

\begin{remark}\label{rmk: leibniz rule comm}
     Since 
     the commutator satisfies the Leibniz rule, for any $u_1,u_2 \in U(L)$ we have:
    \begin{enumerate}
        \item $[x_1,x_2]x_3x_4^{u_{1}}\in \langle  [x_1,x_2]x_3^{u_{1}}\rangle_{T_L} $; 

        \vspace{1ex}
        \item  $ x_1^{u_{2}}x_2 [x_3, x_4]  \in \langle   x_1^{u_{2}}[x_2, x_3]\rangle_{T_L}$; 
        \vspace{1ex}
         \item $  x_1^{u_{1}}x_2[x_3,x_4]x_5^{u_{2}},  x_1^{u_{1}}[x_2,x_3]x_4x_5^{u_{2}} \in \langle x_1^{u_{1}}[x_2,x_3]x_4^{u_{2}} \rangle_{T_L}$.
    \end{enumerate}
\end{remark}

\begin{theorem}\label{teo: Id^L UT_3 e1-e2}
   Let $L$ be a Lie algebra and $UT_3^{\e_1-\e_2}$  be the $L$-algebra $UT_3$ where $L$ acts via the one-dimensional Lie subalgebra of $\D(UT_3)$ spanned by the inner derivation $\e_1-\e_2 =\ad_{e_{22}}$. Then  
there exists a basis $\mathcal{B}_L=\{d_i \mid i \geq 1\}$ of $L$ over $F$ such that 
the $T_L$-ideal of $L$-identities of $UT_3^{\e_1-\e_2}$ is generated by the following $L$-polynomials
    \begin{align*} 
       & x^{d_i}, \quad x^{d_1^3}-x^{d_1}, \quad  x_1^{u_{23}}x_2^{u_{12}}, \quad [x_1,x_2]x_3^{u_{12}}, \quad  x_1^{u_{23}}[x_2,x_3],\quad  x_1^{u_{12}} \big([x_2,x_3]^{u_{23}}-[x_2,x_3]\big), \\  & \big([x_1,x_2]^{u_{12}}-[x_1,x_2]\big) x_3^{u_{23}},\quad \big([x_1,x_2]^{u_{12}}-[x_1,x_2]\big)[x_3,x_4],
    \end{align*}
for all $i\geq 2,$ where $u_{12}:=2^{-1}(d_1^2+d_1)$ and $u_{23}:=2^{-1}(d_1^2-d_1).$ 

Moreover, $c_n^L(UT_3^{\e_1 -\e_2})=(n^2-n)3^{n-2}+(3n-2)2^{n-1}+2$.
\end{theorem}

\begin{proof}
We follow the scheme of Remark~\ref{rmk: general-strategy}, with $\theta(d_1)=\e_1-\e_2$; recall
that $\e_1-\e_2$ acts on the matrix units as in \eqref{eq: nu_12}--\eqref{eq: nu_23}, via
$u_{12}=2^{-1}(d_1^2+d_1)$ and $u_{23}=2^{-1}(d_1^2-d_1)$, whose action annihilates every matrix
unit except, respectively, $e_{12}$ and $e_{23}$.

\textbf{Step 2} is immediate from \eqref{eq: nu_12}--\eqref{eq: nu_23}.

\textbf{Step 3.} Let $f\in P_n^L$. As in the proof of Theorem~\ref{Teo: Id^L UT_3 epsilon1}, the exponents occurring
modulo $I$ reduce to $\{1_{U(L)}, d_1, d_1^2\}$ and as a consequence to $\{1_{U(L)},u_{12},u_{23}\}$, at most two left-normed commutators occur
(Lemma~\ref{lem: conseq 3 commutators 1}), and, by Lemmas~\ref{lem: cons IdUT3e1-e2 1}
and~\ref{lem: cons IdUT3e1-e2 3}, at most one variable may carry exponent $u_{12}$ and at most one
may carry exponent $u_{23}$; so a monomial contains $0$, $1$, or $2$ variables with exponent
different from $1_{U(L)}$.

If none occurs, the argument is exactly as in Theorem~\ref{Teo: Id^L UT_3 ordinarie} (the trivial
action case), giving the ordinary family
$$
x_{i_1}\cdots x_{i_r}[x_{j_1},\ldots,x_{j_s}][x_{k_1},\ldots,x_{k_t}], \qquad
j_1>j_2<\cdots<j_s, \quad k_1>k_2<\cdots<k_t.
$$
If exactly one occurs, with exponent $u_{12}$ (resp.\
$u_{23}$), following the argument as in the proof of Theorem~\ref{Teo: Id^L UT_3 epsilon1} for the
single exponent $d_1$, using \eqref{eq: tail-to-commutator} (applied to $u_{12}$, resp.\ $u_{23}$) in
place of the analogous identity there, and Lemma~\ref{lem: cons IdUT3e1-e2 5} (resp.\
Lemma~\ref{lem: cons IdUT3e1-e2 6})  in place of the Leibniz rule argument fixing the order of the commutator in which appears $u_{12}$ (resp.\ $u_{23}$); this gives the families 
\begin{align*}
&x_{i_{1}}\dots x_{i_{r}}[x_{h_1}^{u_{12}},x_{h_2},\dots,x_{h_{q_1}}][x_{m_1},x_{m_2},\dots,x_{m_{w_1}}],\\ 
&x_{i_{1}}\dots x_{i_{r}}[x_{m'_1},x_{m'_2},\dots,x_{m'_{w_2}}][x_{h'_1}^{u_{23}},x_{h'_2},\dots,x_{h'_{q_2}}],
\end{align*}
where $r,w_1,w_2\geq 0$,  $w_1,w_2\neq 1$, $ q_1,q_2\geq 1,$  $i_1<\cdots < i_r$, $m_1>m_2< \cdots < m_{w_1},$  $m'_1>m'_2< \cdots < m'_{w_2},$  $h_2< \cdots <h_{q_1}$  $h'_2< \cdots <h'_{q_2}.$ Additionally,  if $w_1\geq 2$  then $h_1<h_2$   and if $w_2\geq 2$  then $h'_1<h'_2.$
    

The genuinely new case is when both $u_{12}$ and $u_{23}$ occur together in one monomial. Since
$x_1^{u_{23}}x_2^{u_{12}}\in I$ and, by Lemma~\ref{lem: cons IdUT3e1-e2 3}, also
$x_1^{u_{23}}x_2x_3^{u_{12}}\in I$, the two corresponding variables must appear, modulo $I$, in the fixed
order $u_{12}$ before $u_{23}$.  We examine the cases where the number of commutators $b$ is equal to $2,1,0$ in turn.



\smallskip
\noindent\emph{Case $b=2$.} Since $x_1^{u_{23}}[x_2,x_3]\in I$ and its Leibniz consequence
(Remark~\ref{rmk: leibniz rule comm}), $u_{23}$ must lie in the second commutator $c_2$; since
$$
x_1^{u_{12}}[x_2,x_3]x_4^{u_{23}}\in I \qquad \text{(Lemma~\ref{lem: cons IdUT3e1-e2 4})}
$$
and its Leibniz consequences (Remark~\ref{rmk: leibniz rule comm}), $u_{12}$ must lie in the first commutator $c_1$;
both are placed in the first entry of their commutator.

\smallskip


\noindent\emph{Case $b=1$.} As above, $u_{23}$ must lie in the single commutator, while $u_{12}$
may initially lie in the tail or in the commutator. In the former case, applying \eqref{eq: tail-to-commutator} to $u_{12}$ rewrites the polynomial as a linear combination of terms with two separate commutators, one containing $u_{12}$ and the other containing $u_{23}$. In the latter case, Lemma~\ref{lem: cons IdUT3e1-e2 4} ensures that $u_{12}$ and $u_{23}$ occupy the first two positions of the commutator.  Since $x_1^{u_{23}}x_2^{u_{12}}\in I$, we have
$$
[x_1^{u_{12}},x_2^{u_{23}},x_3]\equiv[x_1^{u_{12}}x_2^{u_{23}},x_3]\pmod I,
$$
and Lemma~\ref{lem: x_1^{u_1}x_2^{u_2}} splits this into two separate commutators, one for each
exponent.

\smallskip
\noindent\emph{Case $b=0$.} Applying \eqref{eq: tail-to-commutator} first moves the variable with
exponent $u_{23}$ into a commutator, giving a linear combination of terms
$$
x_{i_1}\cdots x_{i_{a-1}}x_{i_a}^{u_{12}}x_{i_{a+1}}\cdots x_{i_l}
[x_{j_1}^{u_{23}},x_{j_2},\ldots,x_{j_{n-l}}];
$$
applying again \eqref{eq: tail-to-commutator} to the remaining tail variable with exponent $u_{12}$ then
yields a linear combination of terms
$$
x_{i_1}\cdots x_{i_r}[x_{k_1}^{u_{12}},x_{k_2},\ldots,x_{k_m}]
[x_{j_1}^{u_{23}},x_{j_2},\ldots,x_{j_{n-r-m}}].
$$

\smallskip
In each of these three sub-cases, modulo $I$, both  variables with exponent different from $1_{U(L)}$ lie in the first entries of
two separate commutators, with $u_{12}$ leading. Ordering the remaining entries via
Lemma~\ref{lem: ordered comm x_1^u}, and fixing $l_1<l_2$, $l_1'<l_2'$ via
Lemmas~\ref{lem: cons IdUT3e1-e2 5} and~\ref{lem: cons IdUT3e1-e2 6}, we obtain, modulo $I$, the
family $x_{i_1}\cdots x_{i_r}[x_{l_1}^{u_{12}},\ldots,x_{l_{v_1}}][x_{l_1'}^{u_{23}},\ldots,x_{l_{v_2}'}]$.

Combining the four families above, we conclude that the following $L$-polynomials generate
$P_n^L$ modulo $P_n^L\cap I$:
\begin{equation}\label{No IdUT_3^e_1-e_2}
\begin{split}
&x_{i_1}\dots x_{i_r}[x_{j_1},x_{j_2},\dots,x_{j_s}][x_{k_1},x_{k_2},\dots,x_{k_t}],\\
&x_{i_1}\dots x_{i_r}[x_{h_1}^{u_{12}},x_{h_2},\dots,x_{h_{q_1}}][x_{m_1},x_{m_2},\dots,x_{m_{w_1}}],\\
&x_{i_1}\dots x_{i_r}[x_{m'_1},x_{m'_2},\dots,x_{m'_{w_2}}][x_{h'_1}^{u_{23}},x_{h'_2},\dots,x_{h'_{q_2}}],\\
&x_{i_1}\dots x_{i_r}[x_{l_1}^{u_{12}},x_{l_2},\dots,x_{l_{v_1}}][x_{l'_1}^{u_{23}},x_{l'_2},\dots,x_{l'_{v_2}}],
\end{split}
\end{equation}
where $r,s,t,w_1,w_2\ge0$, with $s,t,w_1,w_2\ne1$, and $q_1,q_2,v_1,v_2\ge1$ such that
$i_1<\cdots<i_r$, $j_1>j_2<\cdots<j_s$, $k_1>k_2<\cdots<k_t$, $h_2<\cdots<h_{q_1}$,
$m_1>m_2<\cdots<m_{w_1}$, $h'_2<\cdots<h'_{q_2}$, $m'_1>m'_2<\cdots<m'_{w_2}$,
$l_1<\cdots<l_{v_1}$, $l'_1<\cdots<l'_{v_2}$. Additionally, if $w_1\ge2$, then $h_1<h_2$, and if
$w_2\ge2$, then $h'_1<h'_2$.


\textbf{Step 4} is as in the proof of Theorem~\ref{Teo: Id^L UT_3 epsilon1}: since $UT_3$ has a
unit, we may, without loss of generality, disregard the tail $x_{i_1}\cdots x_{i_r}$. Hence, if  $f\in \I^L(UT_3^{\e_1-\e_2})$ is a linear combination of the polynomials in \eqref{No IdUT_3^e_1-e_2},   we may
suppose that
\begin{align*}
    f = &\sum_{\mathcal{J}_{j_1}, \mathcal{K}_{k_1}} \alpha_{ \mathcal{J}_{j_1}, \mathcal{K}_{k_1}} [x_{j_1},x_{j_2},\dots,x_{j_s}][x_{k_1},x_{k_2},\dots,x_{k_t}]\\
    &+ \sum_{ \mathcal{H}_{h_1}, \mathcal{M}_{m_1}}  \beta_{\mathcal{H}_{h_1}, \mathcal{M}_{m_1}} [x_{h_1}^{u_{12}},x_{h_2},\dots,x_{h_{q_1}}][x_{m_1},x_{m_2},\dots,x_{m_{w_1}}]\\
    &+ \sum_{ \mathcal{M}'_{m'_1}, \mathcal{H}'_{h'_1}}  \gamma_{ \mathcal{M}'_{m'_1}, \mathcal{H}'_{h'_1}}[x_{m'_1},x_{m'_2},\dots,x_{m'_{w_2}}][x_{h'_1}^{u_{23}},x_{h'_2},\dots,x_{h'_{q_2}}]\\
    &+ \sum_{\mathcal{L}, \mathcal{L}'}  \lambda_{\mathcal{L}, \mathcal{L}'}[x_{l_1}^{u_{12}},x_{l_2},\dots,x_{l_{v_1}}][x_{l'_1}^{u_{23}},x_{l'_2},\dots,x_{l'_{v_2}}],
\end{align*}
where where $\mathcal{J}_{j_1}=\{j_1,\ldots,j_s\},$ $\mathcal{K}_{k_1}=\{k_1,\ldots,k_t\},$ $\mathcal{H}_{h_1}=\{h_1,\ldots,h_{q_1}\},$ $\mathcal{M}_{m_1}=\{m_1,\ldots,m_{w_1}\},$ $\mathcal{H}'_{h'_1}=\{h'_1,\ldots,h'_{q_2}\},$ $\mathcal{M}'_{m'_1}=\{m'_1,\ldots,m'_{w_2}\},$ $\mathcal{L}=\{l_1,\ldots,l_{v_1}\}$ and $\mathcal{L}'=\{l'_1,\ldots,l'_{v_2}\}$
 are subsets of indices in $\{1,\ldots, n\}$ such that $\mathcal{J}_{j_1}\cap \mathcal{K}_{k_1}=\mathcal{H}_{h_1}\cap \mathcal{M}_{m_1}=\mathcal{H}'_{h'_1}\cap \mathcal{M}'_{m'_1}=\mathcal{L}\cap \mathcal{L}'=\emptyset$, with
$s\geq 2,$ $t, w_1,w_2\geq 0,$ $t,w_1,w_2\neq 1,$ $q_1,q_2,v_1,v_2\geq 1,$ $s+t=q_1+w_1=w_2+q_2=v_1+v_2=n,$ and subjected  to the conditions $j_1>j_2<  \cdots < j_s,$ $k_1>k_2<  \cdots < k_t,$ $m_1> m_2< \cdots < m_{w_1},$ $m'_1> m'_2< \cdots < m'_{w_2},$
$l_1< \cdots < l_{v_1},$ $l'_1< \cdots < l'_{v_2},$ $h_2< \cdots <h_{q_1},$ $h'_2< \cdots <h'_{q_2}.$ Additionally, if $w_1\geq 2,$ then $h_1< h_2$, and  if $w_2\geq 2,$ then $h'_1< h'_2.$



\textbf{Step 5.} We show that the $L$-polynomials \eqref{No IdUT_3^e_1-e_2} are linearly independent modulo
$\I^L(UT_3^{\e_1-\e_2})$; by \eqref{eq: nu_12} and \eqref{eq: nu_23}, we make suitable evaluations
to prove that all the coefficients of $f$ are zero.

Let $\mathcal J_{j_1}=\{j_1,\ldots,j_n\}$ be a fixed set such that $j_1>j_2<\cdots<j_n$. If we
evaluate $x_{j_1}=e_{13}$ and $x_{j_2}=\cdots=x_{j_n}=e_{33}$, then we get
$\alpha_{\mathcal J_{j_1},\emptyset}\,e_{13}=0$, and hence $\alpha_{\mathcal J_{j_1},\emptyset}=0$.

Now, let $\mathcal H_{h_1}=\{h_1,\ldots,h_n\}$ be a fixed set such that $h_2<\cdots<h_n$. By
evaluating $x_{h_1}=e_{12}$ and $x_{h_2}=\cdots=x_{h_n}=e_{22}$, we get
$\beta_{\mathcal H_{h_1},\emptyset}\,e_{12}=0$, and so $\beta_{\mathcal H_{h_1},\emptyset}=0$.
Also, by making the evaluation $x_{h_1}=e_{23}$ and $x_{h_2}=\cdots=x_{h_n}=e_{33}$, we get
$\gamma_{\emptyset,\mathcal H_{h_1}}\,e_{23}=0$, and hence $\gamma_{\emptyset,\mathcal H_{h_1}}=0$.

Next, let $\mathcal J_{j_1}=\{j_1,\ldots,j_s\}$ and $\mathcal K_{k_1}=\{k_1,\ldots,k_t\}$ be two
fixed disjoint subsets of $\{1,\ldots,n\}$ such that $s,t\ge2$, $s+t=n$ and
$j_1>j_2<\cdots<j_s$, $k_1>k_2<\cdots<k_t$. If we evaluate $x_{j_1}=e_{12}$,
$x_{j_2}=\cdots=x_{j_s}=e_{11}$, $x_{k_1}=e_{23}$ and $x_{k_2}=\cdots=x_{k_t}=e_{33}$, then we get
$\alpha_{\mathcal J_{j_1},\mathcal K_{k_1}}\,e_{13}=0$ and hence $\alpha_{\mathcal J_{j_1},
\mathcal K_{k_1}}=0$. Here the polynomials of the form
$[x_{h_1}^{u_{12}},x_{h_2},\ldots,x_{h_s}][x_{k_1},x_{k_2},\ldots,x_{k_t}]$ evaluate to zero since
$j_1>j_2<\cdots<j_s$ whereas $h_1<h_2<\cdots<h_s$ and $e_{11}^{\nu_{12}}=0$. Similarly, the
polynomials of the form
$[x_{j_1},x_{j_2},\ldots,x_{j_s}][x_{h_1'}^{u_{23}},x_{h_2'},\ldots,x_{h_t'}]$ evaluate to zero
since $k_1>k_2<\cdots<k_t$ whereas $h_1'<h_2'<\cdots<h_t'$ and $e_{33}^{\nu_{23}}=0$. Also, with
the same reasoning, all the polynomials of the type
$[x_{l_1}^{u_{12}},x_{l_2},\ldots,x_{l_{v_1}}][x_{l_1'}^{u_{23}},x_{l_2'},\ldots,x_{l_{v_2}'}]$
evaluate to zero.

Let $\mathcal H_{h_1}=\{h_1,\ldots,h_{q_1}\}$ and $\mathcal M_{m_1}=\{m_1,\ldots,m_{w_1}\}$ be two
fixed disjoint subsets of $\{1,\ldots,n\}$ with $q_1\ge1$, $w_1\ge2$, $q_1+w_1=n$ and subject to
the conditions $h_1<h_2<\cdots<h_{q_1}$, $m_1>m_2<\cdots<m_{w_1}$. If we make the evaluation
$x_{h_1}=e_{12}$, $x_{h_2}=\cdots=x_{h_{q_1}}=e_{11}$, $x_{m_1}=e_{23}$ and
$x_{m_2}=\cdots=x_{m_{w_1}}=e_{33}$, we get $\beta_{\mathcal H_{h_1},\mathcal M_{m_1}}\,e_{13}=0$,
that is $\beta_{\mathcal H_{h_1},\mathcal M_{m_1}}=0$.

Similarly, if we evaluate $x_{h_1}=e_{23}$, $x_{h_2}=\cdots=x_{h_{q_1}}=e_{33}$, $x_{m_1}=e_{12}$
and $x_{m_2}=\cdots=x_{m_{w_1}}=e_{11}$, we get $\gamma_{\mathcal M_{m_1},\mathcal
H_{h_1}}\,e_{13}=0$, and so $\gamma_{\mathcal M_{m_1},\mathcal H_{h_1}}=0$.

Finally, fixed $\mathcal L=\{l_1,\ldots,l_{v_1}\}$ and $\mathcal L'=\{l_1',\ldots,l_{v_2}'\}$
disjoint subsets of $\{1,\ldots,n\}$ such that $v_1,v_2\ge1$, $v_1+v_2=n$ and $l_1<\cdots<l_{v_1}$,
$l_1'<\cdots<l_{v_2}'$, then from the evaluation $x_{l_1}=e_{12}$,
$x_{l_2}=\cdots=x_{l_{v_1}}=e_{11}$, $x_{l_1'}=e_{23}$ and $x_{l_2'}=\cdots=x_{l_{v_2}'}=e_{33}$,
it follows that $\lambda_{\mathcal L,\mathcal L'}\,e_{13}=0$ and consequently
$\lambda_{\mathcal L,\mathcal L'}=0$.

Therefore, all the scalars appearing in $f$ are zero, and consequently, the $L$-polynomials
\eqref{No IdUT_3^e_1-e_2} are linearly independent modulo $\I^L(UT_3^{\e_1-\e_2})$. This proves that
$\I^L(UT_3^{\e_1-\e_2})=I$ and the $L$-polynomials \eqref{No IdUT_3^e_1-e_2} are a basis of $P_n^L$
modulo $P_n^L\cap\I^L(UT_3^{\e_1-\e_2})$.

  
  Next we determine the $L$-codimensions of $UT_3^{\e_1-\e_2}$ by counting the $L$-polynomials \eqref{No IdUT_3^e_1-e_2}. 
Hence, by Theorems \ref{Teo: Id^L UT_3 ordinarie} and \ref{Teo: Id^L UT_3 epsilon1}, we get
\begin{align*}
     c_n^L(UT_3^{\e_1-\e_2}) =&   2 c_n^L(UT_3^{\e_1})-c_n^L(UT_3)+\sum_{r=0}^{n-2} \binom{n}{r} \left( \sum_{v=1}^{n-r-1}\binom{n-r}{v}\right)\\
      = &2(n^2-4n)3^{n-2}+ 2(3n-2)2^{n-1} +4-(n^2-7n+9) 3^{n-2}-(3n-6)2^{n-1}-3\\
     &+3^n-2^{n+1}+1\\
      = &(n^2 - n)3^{n-2} + (3n - 2)2^{n-1} + 2,
\end{align*}
as claimed.
\end{proof}

\end{document}
